\needspace{12\baselineskip}
\SourceLeaf{319}
\pdfbookmark[1]{16. Über die Zodiaken der Himmelskörper}{zusatz-016}
\begin{center}
[Über die Zodiaken der Himmelskörper.]
\end{center}

\begin{center}
[16.]
\end{center}

\par\noindent\rule{\textwidth}{.3pt}\par

\noindent\small Astronomische Nachrichten. Band XXVII. Nr. 625. Seite 1. 29. Januar 1848.\normalsize

\par\noindent\rule{\textwidth}{.3pt}\par

\begin{center}
Schreiben des Herrn Geheimen Hofrats \textsc{Gauss} an den Herausgeber.

Göttingen, 23. November 1847.
\end{center}

Der in Nr. 615 der A.\,N. abgedruckte zweite Satz von Bahnelementen der Iris, den Herr Prof. \textsc{Goldschmidt} gleich nach dem 19\textsuperscript{.} September berechnet hatte, stimmte auch in den beiden folgenden Monaten noch gut mit den Beobachtungen überein (am 17\textsuperscript{.} und 18\textsuperscript{.} November betrug die Differenz zu Herrn \textsc{Rümkers} Meridianbeobachtung \(16^{\prime\prime}\) in Rektaszension\ und \(2^{\prime\prime}\) in Deklination). Man kann diese Bahnelemente daher schon als sehr gute Näherung betrachten, und ich habe deshalb Herrn Prof. \textsc{Goldschmidt} veranlasst, daraus den Zodiakus der Iris zu berechnen. Das Ergebnis dieser Arbeit gebe ich nachstehend wieder, da es die Suche nach demselben Objekt in früheren Beobachtungen erleichtern kann.

Bei dieser Gelegenheit nehme ich den Aufsatz\HistoricalNoteEight{Band VI, S. 106.]} wieder zur Hand, in dem ich vor 44 Jahren das allgemeine Verfahren zur Bestimmung der Grenzen eines solchen Zodiakus angegeben habe. Ich sehe, dass ich darin auch schon die Ausnahmefälle angedeutet habe, in denen der Bereich der Himmelskugel, in dem ein nach den \textsc{Kepler}schen Gesetzen um die Sonne kreisender Himmelskörper geozentrisch erscheinen kann, entweder gar keine oder nur eine Grenze hat, obwohl das Verfahren im ersten Fall eine geschlossene
\SourceLeaf{320}Kurve, im zweiten zwei solche Kurven ergibt. Dort wird auch die Frage aufgeworfen, was die durch Rechnung gefundenen Kurven in solchen Ausnahmefällen eigentlich bedeuten. Ich habe mich damals auf diese Andeutungen beschränkt, weil eine ausführlichere Behandlung dort ein Horsd’oeuvre gewesen wäre und ich auch gern anderen das Vergnügen lassen wollte, sich mit einer meiner Meinung nach nicht uninteressanten mathematischen Aufgabe zu beschäftigen. Da ich jedoch nicht erfahren habe, dass jemand anderes in der langen Zwischenzeit die mir inzwischen völlig entfallene Frage aufgegriffen hätte, nutze ich diese Gelegenheit, um hier wenigstens das Wesentliche zur Beantwortung mitzuteilen.

Ein geozentrischer Ort des betreffenden Planeten (oder Kometen) ergibt sich, wenn man einen Punkt seiner Bahn mit einem Punkt der Erdbahn kombiniert. Derselbe geozentrische Ort kann sich aber auch aus zwei, drei oder vier verschiedenen Kombinationen ergeben. Zur Veranschaulichung lege man durch die Sonne eine Ebene, die senkrecht zu der Geraden steht, die einen vorgegebenen geozentrischen Ort auf der Himmelskugel darstellt, und projiziere sowohl die Erdbahn als auch die Planetenbahn orthografisch auf diese Ebene. Beide Projektionen sind Ellipsen oder, allgemeiner, Kegelschnitte, die sich entweder gar nicht, zweimal oder viermal schneiden. Eine Berührung ist dabei als das Zusammenfallen zweier Schnittpunkte anzusehen. Jeder Schnittpunkt stellt einen geozentrischen Ort dar, der entweder mit dem vorgegebenen Ort übereinstimmt oder ihm auf der Himmelskugel gegenüberliegt. Das hängt davon ab, welcher der beiden Punkte der Planetenbahn und der Erdbahn, deren Projektionen zusammenfallen, höher liegt. Als Oberseite der Projektionsebene gilt dabei diejenige, auf der der vorgegebene geozentrische Ort liegt.

Daraus ist klar, dass ein beliebiger Punkt der Himmelskugel, der als möglicher geozentrischer Ort vorgegeben wird, dies entweder auf gar keine Weise, auf eine Weise oder auf zwei, drei oder vier Weisen sein kann. Für einen gegebenen Planeten gliedert sich so die ganze Oberfläche der Himmelskugel in verschiedene Teile. Die nach meinem Verfahren bestimmten Kurven sind allgemein gesprochen nichts anderes als Trennlinien zwischen zwei Flächenteilen der Kugel, wobei die Punkte auf der einen Seite auf zwei Weisen mehr geozentrische Orte sein können als die auf der
\SourceLeaf{321}anderen Seite. Die Punkte auf den Trennlinien bilden den Übergang, d. h. sie können auf eine Weise mehr als die Punkte auf der einen Seite und auf eine Weise weniger als die Punkte auf der anderen Seite geozentrische Orte sein. Es lässt sich auch ein Kriterium angeben, nach dem man a priori entscheiden kann, auf welcher Seite der Trennlinie es zwei Lösungen mehr gibt als auf der anderen. Darauf möchte ich hier nicht weiter eingehen\HistoricalNoteEight{Siehe Notiz 17.]}.

Alle bisher bekannten Planeten haben Bahnen, bei denen die aus der Theorie gewonnenen Trennlinien stets tatsächliche Grenzen sind. Es gibt nämlich zwei Trennlinien, die die Himmelskugel in drei Flächenbereiche unterteilen. Zwei davon sind voneinander getrennte Flächen, in die gar keine geozentrischen Orte fallen. Der dritte, gürtelförmige Bereich zwischen ihnen enthält alle geozentrischen Orte. Jeder Punkt innerhalb des Gürtels kann auf zwei Weisen und jeder Punkt auf der Grenze auf eine Weise ein geozentrischer Ort sein.

Man kann sich aber hypothetische Bahnen vorstellen\HistoricalNoteDouble{Siehe Notiz 18.]} (mehrere Kometenbahnen könnten vielleicht unter diesen Fall fallen, was ich bisher jedoch nicht untersucht habe), bei denen zwei Grenzen die Himmelskugel ebenfalls in drei Teile teilen. Der eine Teil enthält gar keine geozentrischen Orte, der andere die geozentrischen Orte, die sich auf zwei Weisen ergeben, der dritte hingegen die Punkte, die auf vier verschiedene Weisen geozentrische Orte sein können.

Weitere Vielfalt ergibt sich, wenn eine Trennlinie sich selbst ein- oder zweimal schneidet und so eine einfache oder doppelte Schleife bildet. Im ersten Fall trennt sie zwei, im zweiten drei Flächenbereiche vom gürtelförmigen Teil ab, in denen es entsprechend 4 und 0 oder 4, 0, 4 Lösungen gibt.

Wieder anders verhält es sich mit einer Bahn, die wie ein Kettenglied in die Erdbahn eingreift. In einem solchen Fall gibt es nur eine zusammenhängende Grenzlinie und damit zwei getrennte Flächenteile. Die Punkte des einen Teils können auf eine Weise, die des anderen Teils auf drei Weisen geozentrische Orte sein.

\needspace{29\baselineskip}
\SourceLeaf{322}
\begin{center}
Zodiakus der Iris.
\end{center}
\begingroup\small
\setlength{\tabcolsep}{7pt}
\renewcommand{\arraystretch}{1.15}
\begin{center}
\begin{tabular}{|r|r@{\;}r|r@{\;}r|}
\multicolumn{1}{|c|}{Rektas-}&\multicolumn{4}{c|}{Deklination}\\
\multicolumn{1}{|c|}{zension}&\multicolumn{2}{c|}{nördl. Grenze}&\multicolumn{2}{c|}{südl. Grenze}\\
\(0^{0}\)&\(+12^{0}\)&\(9^{\prime}\)&\(+3^{0}\)&\(44^{\prime}\)\\
5&14&10&5&41\\
10&16&5&7&33\\
15&17&51&9&19\\
20&19&30&10&57\\
25&21&0&12&29\\
30&22&21&13&52\\
35&23&33&15&7\\
40&24&36&16&14\\
45&25&29&17&11\\
50&26&12&17&59\\
55&26&45&18&39\\
60&27&9&19&8\\
65&27&24&19&29\\
70&27&29&19&41\\
75&27&24&19&42\\
80&27&11&19&36\\
85&26&48&19&19\\
90&26&15&18&54\\
95&25&33&18&19\\
100&24&41&17&35\\
105&23&41&16&42\\
110&22&32&15&41\\
115&21&14&14&31\\
120&19&47&13&12\\\hline
\end{tabular}
\end{center}
\begin{center}
\begin{tabular}{|r|r@{\;}r|r@{\;}r|}
\multicolumn{1}{|c|}{Rektas-}&\multicolumn{4}{c|}{Deklination}\\
\multicolumn{1}{|c|}{zension}&\multicolumn{2}{c|}{nördl. Grenze}&\multicolumn{2}{c|}{südl. Grenze}\\
\(125^{0}\)&\(+18^{0}\)&\(12^{\prime}\)&\(+11^{0}\)&\(46^{\prime}\)\\
130&16&29&10&12\\
135&14&40&8&31\\
140&12&45&6&45\\
145&10&44&4&53\\
150&8&39&2&56\\
155&6&31&\(+0\)&56\\
160&4&22&\(-1\)&6\\
165&\(+2\)&12&3&8\\
170&0&0&5&11\\
175&\(-2\)&9&7&13\\
180&4&14&9&11\\
185&6&16&11&6\\
190&8&12&12&56\\
195&10&2&14&40\\
200&11&45&16&19\\
205&13&20&17&50\\
210&14&46&19&13\\
215&16&4&20&58\\
220&17&13&21&35\\
225&18&12&22&33\\
230&19&1&23&23\\
235&19&40&24&4\\
240&20&10&24&35\\\hline
\end{tabular}
\end{center}
\begin{center}
\begin{tabular}{|r|r@{\;}r|r@{\;}r|}
\multicolumn{1}{|c|}{Rektas-}&\multicolumn{4}{c|}{Deklination}\\
\multicolumn{1}{|c|}{zension}&\multicolumn{2}{c|}{nördl. Grenze}&\multicolumn{2}{c|}{südl. Grenze}\\
\(245^{0}\)&\(-20^{0}\)&\(29^{\prime}\)&\(-24^{0}\)&\(58^{\prime}\)\\
250&20&38&25&11\\
255&20&38&25&15\\
260&20&27&25&11\\
265&20&5&24&56\\
270&19&34&24&33\\
275&18&52&24&1\\
280&18&0&23&19\\
285&16&58&22&28\\
290&15&46&21&28\\
295&14&25&20&20\\
300&12&54&19&2\\
305&11&14&17&36\\
310&9&26&16&2\\
315&7&30&14&21\\
320&5&28&12&34\\
325&3&21&10&40\\
330&\(-1\)&9&8&41\\
335&\(+1\)&6&6&39\\
340&3&22&4&34\\
345&5&37&2&28\\
350&7&51&\(-0\)&22\\
355&10&2&\(+1\)&42\\
360&12&9&3&44\\\hline
\end{tabular}
\end{center}
\endgroup
\par\medskip\begin{center}\rule{.22\textwidth}{.3pt}\end{center}
